% Part: counterfactuals
% Chapter: introduction
% Section: counterfactuals

\documentclass[../../../include/open-logic-section]{subfiles}

\begin{document}

\olfileid{cnt}{int}{cnt}

\olsection{د واقعيت خلاف شرطونه}

په انګرېزۍ کښې د «که \dots نو \dots» جوړښتونو يوه ډېره عامه او
مهمه بڼه د \emph{to be} («وي/وے») فعل د تېر التزامي صورت په کارولو جوړېږي:
«که داسې وے چې \dots نو داسې به وے چې \dots»۔ څرنګه چې د داسې
شرطي مقدم عموماً کاذب، يعنې له واقعيت خلاف وي، دې ته
\emph{د واقعيت خلاف شرطونه} وايي (او څرنګه چې د \emph{to be}
فعل التزامي صورت کاروي، \emph{التزامي شرطونه} هم ورته وايي)۔ د نسخې
څرګندونه (OLCNT-005): په منجمده سرچينه کښې د دې قوس تړون نشت دے؛
دلته تړل شوے دے۔ دا له \emph{خبري} شرطونو بېل دي، چې بڼه يې داسې
ده: «که داسې ده چې \dots نو داسې ده چې \dots»۔ د واقعيت خلاف او
خبري شرطونه په صدق-شرطونو کښې توپير لري۔ د اېډمز مشهوره بېلګه وګورئ:
\begin{quote}
  که اوسوالډ کېنېډي نۀ وي وژلے، نو بل چا وژلے دے۔
  
  که اوسوالډ کېنېډي نۀ وے وژلے، نو بل چا به وژلے وے۔
\end{quote}
لومړۍ خبري ده، دويمه د واقعيت خلاف شرطي ده۔ لومړۍ په څرګنده
رښتينې ده: مونږ پوهېږو چې ولسمشر جان اېف کېنېډي \emph{چا} وژلے
دے، او که هغه کس (د وارېن د راپور خلاف) لي هاروي اوسوالډ نۀ و، نو
بل چا کېنېډي وژلے دے۔ دويمه بيا بېل څه وايي۔ دا دعوه کوي چې که
اوسوالډ کېنېډي نۀ وے وژلے، يعنې که په ډالاس کښې د ډزو د پېښې
مخنيوے شوے وے يا ناکامه شوې وے، نو تاريخ به تر هغې وروسته داسې
مخ ته تلے وے چې بله د وژنې هڅه به بريالۍ شوې وے۔ د دې د رښتيني
کېدو لپاره بايد ځواکمنو قوتونو د جان اېف کېنېډي د مرګ د يقيني
کولو سازش کړے وے (څرنګه چې د کېنېډي د وژنې د سازش د نظريو ډېر پلويان باور لري)۔

لا بحث روان دے چې آيا \emph{خبري} شرطي د مادي شرطي په وسيله سم
بيانېږي، په تېره بيا دا چې آيا د مادي شرطي پاراډوکسونه په داسې
طريقه «تشريح» کېدے شي چې له دې سره برابره وي چې مادي شرطي د
انګرېزۍ د خبري شرطونو صدق-شرطونه ورکوي۔ خو په دې کښې اختلاف نشت
چې د واقعيت خلاف شرطونه د مادي شرطونو په وسيله سم نۀ شي
نښه‌ليکل کېدے۔ دا ځکه څرګنده ده چې که څه هم د واقعيت خلاف شرطونو
مقدمونه عموماً کاذب وي، هر هغه د واقعيت خلاف شرطي چې کاذب مقدم
لري رښتينې نۀ وي---د بېلګې په توګه، که تاسو د وارېن پر راپور
باور لرئ او د جان اېف کېنېډي د وژلو سازش نۀ و، نو د اېډمز د
واقعيت خلاف شرطي يې يوه بېلګه ده۔

د واقعيت خلاف شرطونه په علّي استدلال کښې مهم رول لري: د هغو د
کارولو يوه څرګنده بېلګه د علّي اړيکو بيانول دي۔ د بېلګې په توګه،
د اورلګيت د لرګي وهل د هغه د بلېدو سبب کېږي، او دا په دې وينا
بيانولے شئ: «که دا د اورلګيت لرګے وهل شوے وے، نو بل شوے به وے»۔
مادي شرطونه، او په عموم کښې خبري شرطونه، د دې د بيانولو لپاره
نۀ شي کارېدے: «د اورلګيت لرګے وهل کېږي $\lif$ د اورلګيت لرګے
بلېږي» هله رښتينې ده چې لرګے هېڅکله ونۀ وهل شي، پرته له دې چې
که وهل شوے وے نو څه به شوي وو۔ تر دې هم بده دا ده چې «د
اورلګيت لرګے وهل کېږي $\lif$ د اورلګيت لرګے د ګلونو په ګېډۍ
بدلېږي» هم هله رښتينې ده چې لرګے هېڅکله ونۀ وهل شي، خو که لرګے
وهل شوے وے نو په يقين سره به د ګلونو په ګېډۍ نۀ وے بدل شوے۔

لا بحث روان دے چې د واقعيت خلاف شرطونو سم منطق په دقيقه توګه
څه دے۔ سټالنکر او لويس د واقعيت خلاف شرطونو يوه اغېزمنه شننه
وړاندې کړه۔ د هغو له مخې، د واقعيت خلاف شرطي «که داسې وے
چې~$S$ نو داسې به وے چې~$T$» هله او يوازې هله رښتينې ده چې $T$
په هغه د واقعيت خلاف حالت («ممکنه نړۍ») کښې رښتينے وي چې د
واقعي نړۍ له حالت سره تر ټولو نژدې وي او~$S$ پکښې رښتينے وي۔
دې ته «هستي‌اړونده» شننه وايي، ځکه چې د ممکنو نړۍو هستي‌پوهنې
ته رجوع کوي۔ نورې شننې شرطي احتمالونه يا د باور د بياکتنې
نظريې کاروي۔ د واقعيت خلاف شرطونو ډېر بېلابېل منطقونه وړانديز
شوي دي۔ آن د واقعيت خلاف شرطونو يو يوازينے د لويس--سټالنکر
منطق هم نشت: که څه هم سټالنکر او لويس د تر ټولو نژدې ممکنو
نړۍو په حواله په ورته طريقو شننې وړاندې کړې، د هغو فرضونه د
معتبرو استنتاجونو بېلابېلې پايلې ورکوي۔

\end{document}
